%! TEX program = xelatex
\documentclass[12pt,a4paper,fontset=none]{ctexart}
\usepackage[top=25mm,bottom=25mm,left=28mm,right=28mm]{geometry}
\usepackage{booktabs,tabularx,array,enumitem,fancyhdr,hyperref}
\setCJKmainfont[
  Path=./fonts/,
  UprightFont=SourceHanSerifCN-Regular.otf,
  BoldFont=SourceHanSerifCN-Bold.otf
]{SourceHanSerifCN}
\setCJKfamilyfont{hei}[
  Path=./fonts/,
  UprightFont=FandolHei-Regular.otf,
  BoldFont=FandolHei-Bold.otf
]{FandolHei}
\setCJKfamilyfont{song}[
  Path=./fonts/,
  UprightFont=SourceHanSerifCN-Regular.otf,
  BoldFont=SourceHanSerifCN-Bold.otf
]{SourceHanSerifCN}
\def\songti{\CJKfamily{song}}
\def\heiti{\CJKfamily{hei}}
\newcolumntype{L}{>{\raggedright\arraybackslash}X}
\pagestyle{fancy}
\fancyhf{}
\fancyhead[R]{\zihao{-5}\songti 队伍控制号：TJMML2026036}
\fancyfoot[C]{\thepage}
\renewcommand{\headrulewidth}{0pt}
\setlength{\headheight}{15pt}
\setlength{\parindent}{2em}
\setlength{\parskip}{0.4em}
\hypersetup{hidelinks,pdfauthor={},pdfsubject={第一届天津市五校数学建模联赛支撑材料}}

\begin{document}
\begin{center}
  {\zihao{2}\bfseries 人工智能工具使用详情\par}
\end{center}

\section*{一、所用AI工具}
\begin{table}[htbp]
\centering
\renewcommand{\arraystretch}{1.25}
\begin{tabularx}{\textwidth}{p{3.3cm}L}
\toprule
项目 & 详情 \\
\midrule
工具名称 & OpenAI Codex \\
版本/型号 & GPT-5 \\
开发机构/公司 & OpenAI \\
使用日期 & 2026年8月17日至2026年8月19日 \\
\bottomrule
\end{tabularx}
\end{table}

\section*{二、具体使用目的和环节}
\begin{enumerate}[leftmargin=2.3em,itemsep=0.55em]
  \item \textbf{论文结构梳理与问题拆解。}
  辅助讨论“数据治理—统计预测—条件情景”的论文结构，并对三个子问题的衔接逻辑和表达层次提供建议。

  \item \textbf{数据处理方案与程序实现。}
  辅助讨论缺失值标记、双侧对数线性插值、时间有序滚动检验及数据泄漏防范；辅助编写和精简指数回归、原尺度Ridge、残差自助法、情景恒等式和OAT敏感性分析的实现代码。

  \item \textbf{文字表达与\LaTeX{}排版。}
  辅助优化部分表述的准确性和连贯性，调整表格、公式、图题、引用、附录代码及页眉页码的\LaTeX{}排版。

  \item \textbf{一致性与可复现性检查。}
  辅助执行程序语法检查、表格数值复算、正文与附录方法对应检查以及PDF编译和版面检查。
\end{enumerate}

\section*{三、人工复核说明}
参赛队员对数据来源、统计口径、缺失值处理、模型设定、程序代码、参数、计算结果、图表及论文结论进行了逐项复核，并作出最终取舍和定稿。AI工具仅用于上述辅助环节，未将未经人工核验的生成内容直接作为数据事实、模型结果或论文结论。

\end{document}
